\subsection{Conjuntos}

\begin{itemize}
    \item $C$: Conjunto de locais candidatos para instalação de novos eletropostos.
    \item $E$: Conjunto de eletropostos concorrentes já existentes na malha urbana.
    \item $P$: Conjunto total de Pontos de Interesse (POIs) -- hospitais, educação, transporte, comércio.
    \item $P_{\text{cob}} \subseteq P$: Subconjunto de POIs já atendidos (cobertos) por pelo menos um eletroposto existente $e \in E$.
    \item $P_{\text{ncob}} = P \setminus P_{\text{cob}}$: Subconjunto de POIs não atendidos, que dependem das decisões de instalação em $C$.
    \item $V(c)$: Conjunto de vizinhos espaciais do candidato $c$ (dentro do raio crítico de influência).
    \item $\mathcal{S}_c$: Conjunto de todas as combinações possíveis de vizinhos de $c$ que podem ser designados simultaneamente ($\mathcal{S}_c \subseteq 2^{V(c)}$).
\end{itemize}

\subsection{Parâmetros}

\subsubsection{Lucro e Atratividade (pré-computados via simulador espacial)}

\begin{itemize}
    \item $A_c$: Atratividade gravitacional base do candidato $c$ (Score derivado dos POIs próximos).
    \item $B$: Barreira de viabilidade (custo fixo mínimo ou atratividade mínima requerida para operação neutra).
    \item $\Pi_c$: Penalidade aplicada à atratividade do candidato $c$ devido à proximidade (sobreposição de áreas de influência) com eletropostos já existentes $e \in E$.
    
    \item $R_c^{\text{próprio}}$: Retorno líquido do eletroposto $c$ se operasse sem novos vizinhos, mas \textbf{sujeito à concorrência existente}. Definido como:
    \begin{equation}
        R_c^{\text{próprio}} = A_c - B - \Pi_c
    \end{equation}
    
    \item $\Delta R_{c,c'}$: Retorno (demanda) desviado do candidato $c$ para seu vizinho $c' \in V(c)$ quando $c'$ é designado. Representa a canibalização espacial entre as novas opções, calculada com base na distância de decaimento pela malha viária.
    
    \item $R_c(S_c)$: Retorno esperado do candidato $c$ quando exatamente o conjunto de vizinhos $S_c \subseteq V(c)$ está designado. Definido como:
    \begin{equation}
    R_c(S_c) = R_c^{\text{próprio}} - \sum_{c' \in S_c} \Delta R_{c,c'}
    \end{equation}
\end{itemize}

\subsubsection{Qualidade de rede}

\begin{itemize}
    \item $Q_{\min}^{\text{cob}}$: Cobertura mínima da totalidade de POIs requerida na cidade (percentual, ex: 85\%).
    \item $r_{\text{cob}}$: Raio de cobertura máximo aceitável pela malha viária para um POI ser considerado atendido (metros).
\end{itemize}

\subsection{Variáveis de decisão}

\begin{itemize}
    \item $x_c \in \{0,1\}$: Indicadora -- designar a instalação de um eletroposto no candidato $c$.
    \item $y_p \in \{0,1\}$: Indicadora -- POI $p \in P_{\text{ncob}}$ passa a ser coberto por pelo menos um \textit{novo} eletroposto designado.
    \item $z_{c,S_c} \in \{0,1\}$: Indicadora -- a configuração exata de vizinhos \textit{novos} designados para o eletroposto $c$ é o conjunto $S_c$.
\end{itemize}

\subsection{Função objetivo}

\begin{equation}
\label{eq:objetivo_linear}
\max \quad Z = \sum_{c \in C} \sum_{S_c \in \mathcal{S}_c} R_c(S_c) \cdot z_{c,S_c}
\end{equation}

\textit{Nota: O objetivo maximiza o retorno líquido global, selecionando a configuração de concorrência real ($z_{c,S_c}$) que a rede terminará enfrentando após as designações.}

\subsection{Restrições}

\subsubsection{Grupo 1: Linearização e Configuração de Vizinhos}

\hspace{0.5cm} \textbf{R1: Atribuição única de configuração}

\begin{align}
\sum_{S_c \in \mathcal{S}_c} z_{c,S_c} &= x_c \quad \forall c \in C \label{eq:restricao_config}
\end{align}

\hspace{0.5cm} \textbf{R2: Consistência lógica de designações}

Se a configuração $S_c$ for a escolhida ($z_{c,S_c} = 1$), os vizinhos correspondentes devem estar obrigatoriamente designados ou desligados:

\begin{align}
x_{c'} &\geq z_{c,S_c} \quad \forall c \in C, \forall S_c \in \mathcal{S}_c, \forall c' \in S_c \label{eq:restricao_consistencia_in} \\
x_{c'} &\leq 1 - z_{c,S_c} \quad \forall c \in C, \forall S_c \in \mathcal{S}_c, \forall c' \in V(c) \setminus S_c \label{eq:restricao_consistencia_out}
\end{align}

\subsubsection{Grupo 2: Qualidade Global de Rede (Considerando Infraestrutura Existente)}

\hspace{0.5cm} \textbf{R3: Indicador de nova cobertura para POIs desassistidos}

\begin{align}
y_p &\leq \sum_{c \in C: \text{dist}(c,p) \leq r_{\text{cob}}} x_c \quad \forall p \in P_{\text{ncob}} \label{eq:cobertura_poi}
\end{align}

Um POI desassistido $p$ só será contabilizado como coberto ($y_p = 1$) se pelo menos um novo candidato em seu raio de alcance for instalado.

\hspace{0.5cm} \textbf{R4: Meta integrada de cobertura mínima}

\begin{align}
|P_{\text{cob}}| + \sum_{p \in P_{\text{ncob}}} y_p &\geq \lceil Q_{\min}^{\text{cob}} \cdot |P| \rceil \label{eq:meta_cobertura}
\end{align}

A quantidade de POIs cobertos gratuitamente pela infraestrutura existente ($|P_{\text{cob}}|$) somada aos POIs resgatados pela nova infraestrutura deve atender ou superar o volume mínimo exigido pelo planejamento urbano.